% Part: counterfactuals
% Chapter: minimal-change-semantics
% Section: sphere-models

\documentclass[../../../include/open-logic-section]{subfiles}

\begin{document}

\olfileid{cnt}{min}{sph} % BN-SRC-423: প্রতিবাস্তব অংশের সঠিক file-id।

\olsection{গোলক মডেল}

প্রতিবাস্তব শর্তবচনের আনুষ্ঠানিক অর্থতত্ত্ব দেওয়ার একটি উপায় হল
লুইসের অনানুষ্ঠানিক ব্যাখ্যাকে গাণিতিক কাঠামোয় রূপ দেওয়া। তখন
জগৎ~$w$-কে ঘিরে থাকা গোলকগুলি হবে জগৎসমূহের সেট। গোলকগুলি
পরস্পর-অন্তর্ভুক্ত বলে~$w$-কে ঘিরে জগৎসমূহের সেটগুলিকে উপসেট
সম্পর্ক অনুসারে রৈখিক ক্রমে সাজানো থাকতে হবে।

\begin{defn}
  একটি \emph{গোলক মডেল} হল ত্রয়ী~$\mModel{M} = \tuple{W, O, V}$,
  যেখানে $W$ জগৎসমূহের একটি অশূন্য সেট, $V\colon \PVar \to \Pow{W}$
  একটি মাননির্ধারণ এবং $O\colon W \to \Pow{\Pow{W}}$ প্রতিটি
  জগৎ~$w$-কে একটি \emph{গোলকব্যবস্থা}~$O_w$ নির্ধারণ করে। প্রতিটি
  $w$-এর জন্য $O_w$ হল জগৎসমূহের সেটগুলির একটি সেট এবং তাকে
  নিম্নলিখিত শর্তগুলি পূরণ করতে হবে:
  \begin{enumerate}
  \item $O_w$-এর \emph{কেন্দ্রে} আছে~$w$: $\{w\} \in O_w$।
  \item $O_w$-এর গোলকগুলি \emph{পরস্পর-অন্তর্ভুক্ত}: যখনই $S_1$,
    $S_2 \in O_w$, তখন $S_1 \subseteq S_2$ অথবা $S_2 \subseteq S_1$;
    অর্থাৎ, $\subseteq$ অনুসারে $O_w$ রৈখিক ক্রমে সাজানো।
  \item অশূন্য সেট-সংকলনের সংযুক্তির অধীনে $O_w$ আবদ্ধ।
  \item অশূন্য সেট-সংকলনের ছেদের অধীনে $O_w$ আবদ্ধ।
  \end{enumerate}
\end{defn}

$O_w$-এর পিছনের ধারণা হল, $w$-কে ঘিরে থাকা জগৎগুলি
$w$ থেকে দূরত্ব অনুযায়ী স্তরে সাজানো।
% BN-SRC-766: সংজ্ঞায় খালি গোলক অনুমোদিত; ক্ষুদ্রতম অশূন্য গোলকই নিজের জগৎ।
সংজ্ঞায় শূন্য গোলকও থাকতে পারে; এতে কোনো জগৎ না-থাকায়
পরিতৃপ্তি-শর্তে তার কোনো প্রভাব নেই। একেবারে ভিতরের অশূন্য গোলকটি কেবল $w$
নিজেই, অর্থাৎ সেট~$\{w\}$: অন্য যে-কোনও গোলকের বাকি জগৎগুলির
তুলনায় $w$ নিজে~$w$-এর কাছেই বেশি নিকটবর্তী। যদি $S \subsetneq S'$
হয়, তবে $S' \setminus S$-এর জগৎগুলি~$S$-এর জগৎগুলির চেয়ে
$w$ থেকে দূরে: $S' \setminus S$ হল $S$ ও $S'$-এর বাইরের জগৎগুলির
মাঝের ``স্তর''। বিশেষ করে, প্রতিটি গোলককে তার বাইরের সীমানার
ভিতরে থাকা সব জগৎ-সহ বুঝতে হবে; গোলকগুলি কেবল পৃথক স্তর নয়।

\begin{figure}
\begin{center}
\begin{tikzpicture}[layerwidth=1.5,scale=.6]\tiny
  \spheresystem{3}
  \propositionintersect{0}{3}{60}{6}
  \path (0,0) node[world] {$w$};
  \spherepos{-30}{2}{node[world] {$w_2$}}
  \spherepos{220}{2}{node[world] {$w_3$}}
  \spherepos{90}{2}{node[world] {$w_1$}}
  \spherepos[shift={(.1,0)}]{10}{3}{node[world] {$w_5$}}
  \spherepos[shift={(.1,0)}]{-10}{3}{node[world] {$w_6$}}
  \spherepos{225}{3}{node[world] {$w_4$}}
  \spherepos{20}{4}{node[world] {$w_7$}}
  \path (5,0) node[left] {$p$};
\end{tikzpicture}
\caption{গোলক মডেলের চিত্র}
\ollabel{fig:sphere-model}
\end{center}
\end{figure}

\olref{fig:sphere-model}-এর চিত্রটি সেই গোলক মডেল দেখায়, যেখানে
$W = \{w, w_1, \dots, w_7\}$ এবং $V(p) = \{w_5, w_6,
w_7\}$। সবচেয়ে ভিতরের গোলক $S_1 = \{w\}$। $w$-এর সবচেয়ে
কাছের জগৎগুলি $w_1, w_2, w_3$; তাই পরের বড় গোলকটি
$S_2 = \{w, w_1, w_2, w_3\}$। আরও বাইরের জগৎগুলি হল
$w_4$, $w_5$, $w_6$; তাই সবচেয়ে বাইরের গোলক
$S_3 = \{w, w_1, \dots, w_6\}$। $w$-কে ঘিরে গোলকব্যবস্থাটি
$O_w = \{S_1, S_2, S_3\}$। জগৎ~$w_7$ $w$-কে ঘিরে কোনও গোলকে
নেই। যে নিকটতম জগৎগুলিতে $p$ সত্য, সেগুলি হল $w_5$ ও~$w_6$;
তাই $p$-স্বীকারক ক্ষুদ্রতম গোলকটি হল~$S_3$।

গোলক মডেল~$\mModel M$-এ জগৎ~$w$-তে সূত্র~$!A$-এর পরিতৃপ্তি,
$\mSat{M}{!A}[w]$, সংজ্ঞায়িত করতে বচনমূলক !!{formula}s-এর
% BN-SRC-771: বর্তমান ভাষা বচনমূলক সংযোজক ও প্রতিবাস্তব সংযোজকের; ত্রয়ীতে আলাদা মোডাল R নেই।
সংজ্ঞায় $!B \cif !C$-এর জন্য একটি শর্ত যোগ করি:

\begin{defn}
  $\mSat{M}{!B \cif !C}[w]$ তখন এবং কেবল তখনই, যখন নিচের অন্তত একটি শর্ত পূরণ হয়:
  \begin{enumerate}
  \item\ollabel{sphere-vac} সব $u \in \bigcup O_w$-এর জন্য
    $\mSat/{M}{!B}[u]$; অথবা
  \item\ollabel{sphere-nonvac} এমন কোনও $S \in O_w$ আছে, যার জন্য
    \begin{enumerate}
      \item কোনও $u \in S$-এর জন্য $\mSat{M}{!B}[u]$; এবং
      \item সব $v \in S$-এর জন্য হয় $\mSat/{M}{!B}[v]$, নয়তো
        $\mSat{M}{!C}[v]$।
    \end{enumerate}
  \end{enumerate}
\end{defn}

এই সংজ্ঞা অনুসারে, $\mSat{M}{!B \cif !C}[w]$ তখন এবং কেবল তখনই, যখন হয়
পূর্বপক্ষ~$!B$ $w$-কে ঘিরে সব গোলকের সর্বত্র মিথ্যা, নয়তো
এমন একটি গোলক~$S$ আছে যেখানে $!B$ সত্য এবং সেই
``$!B$-স্বীকারক'' গোলকের প্রতিটি জগতে বস্তুগত শর্তবচন
$!B \lif !C$ সত্য। স্বজ্ঞাগত ব্যাখ্যা থেকে যা মনে হতে পারে,
তার বিপরীতে সংজ্ঞায় $S$-কে \emph{সবচেয়ে ভিতরের}
$!B$-স্বীকারক গোলক হতে বলা হয়নি। কিন্তু কোনও গোলক~$S$-এ
% BN-SRC-422: শুধু ছোট পূর্বপক্ষ-স্বীকারক গোলকগুলি অশূন্য clause সংরক্ষণ করে।
\olref{sphere-nonvac}-এর শর্ত পূরণ হলে, গোলক~$S$-এর অন্তর্গত যত ছোট
পূর্বপক্ষ-স্বীকারক গোলক আছে, সেগুলিতেও শর্তটি পূরণ হয়; অতএব
সবচেয়ে ভিতরের এমন গোলক থাকলে সেখানেও হয়।

% BN-SRC-769: কঠোর অসীম ধারা একাই ক্ষুদ্রতম স্বীকারক গোলকের অনস্তিত্ব প্রমাণ করে না; ধারার criterion-এ ভিতর পর্যন্ত পৌঁছনোর শর্ত চাই।
আরও লক্ষণীয়, গোলক মডেলের সংজ্ঞা সবচেয়ে ভিতরের
$!B$-স্বীকারক গোলক \emph{থাকতেই হবে} এমন দাবি করে না:
$!B$-স্বীকারক গোলকগুলির অসীম ধারা
$S_1 \supsetneq S_2 \supsetneq \dots \supsetneq \{w\}$ থাকতে পারে,
তবে শুধু এই কঠোর ধারা থাকাই ক্ষুদ্রতম স্বীকারক গোলক না-থাকার
প্রমাণ নয়: সব স্বীকারক গোলকের ছেদেও পূর্বপক্ষ সত্য এমন জগৎ
থাকলে সেই ছেদই ক্ষুদ্রতম। সেই ছেদে পূর্বপক্ষ সত্য এমন কোনো
জগৎ না-থাকলে কোনও সবচেয়ে ভিতরের $!B$-স্বীকারক গোলক নেই।
ধারাটির কোনও সদস্য প্রত্যেক পূর্বপক্ষ-স্বীকারক গোলকের মধ্যে
অন্তর্ভুক্ত হয়, এমনভাবে ধারাটি নেওয়া হলে তবেই তার সাহায্যে
দ্বিদিকের নিচের শর্ত বলা যায়: কোনও~$i$ থেকে শুরু করে $S_i$, $S_{i+1}$, \dots
গোলকগুলির সর্বত্র $!B \lif !C$ সত্য হলে এবং কেবল তখনই
$\mSat{M}{!B \cif !C}[w]$।

\end{document}
